\documentclass[11pt]{article}
\usepackage[margin=1in]{geometry}
\usepackage{amsmath,amssymb,amsthm,mathtools}
\usepackage{graphicx}
\usepackage{microtype}
\usepackage{xcolor}
\usepackage{hyperref}
\usepackage{enumitem}

\hypersetup{colorlinks=true,linkcolor=blue!55!black,citecolor=blue!55!black,urlcolor=blue!55!black}
\setlist{nosep,leftmargin=*}

\newtheorem{theorem}{Theorem}
\newtheorem{corollary}[theorem]{Corollary}
\newtheorem{lemma}[theorem]{Lemma}
\newcommand{\Sp}{S_p}
\newcommand{\hp}{\widetilde h_{S_p}}
\newcommand{\E}{\mathbb E}
\newcommand{\N}{\mathbb N}
\newcommand{\Mmap}{\mathcal M}

\title{Sharp Universal Metric Distortion into Every Schatten Class}
\author{Candidate full solution for arXiv:2405.03517, Question 2}
\date{19 July 2026}

\begin{document}
\maketitle

\begin{center}
\fcolorbox{blue!55!black}{blue!4}{\parbox{0.91\textwidth}{
\textbf{Status.} Candidate full solution, likely valid, pending human review.
The argument answers the source question for every $1\le p<\infty$,
including the trace-class endpoint $p=1$.}}
\end{center}

\begin{abstract}
Escudero Gutierrez and Muguruza ask for the worst-case distortion
$D_{n,S_p}$ of an $n$-point metric space into the Schatten class $S_p$.
We prove, with universal implicit constants,
\[
 D_{n,S_p}\asymp \max\left\{1,\frac{\log n}{p}\right\}
 \qquad(1\le p<\infty).
\]
For $p\ge2$, a lazy-walk interpolation argument gives
$\widetilde h_{S_p}(G)\ge\delta/(14p)$ for every regular graph with
normalized adjacency gap $\delta$.  The missing range $1\le p\le2$
requires a different mechanism.  Translate a Schatten-valued map at a
well-chosen vertex and apply the noncommutative Mazur map
$M_{p,2}(x)=u|x|^{p/2}$ into $S_2$.  Ricard's uniform Holder and inverse
Lipschitz estimates, a weighted averaging lemma, and the Hilbert Poincare
inequality yield
\[
 \widetilde h_{S_p}(G)\ge (\delta/K)^{1/p}\qquad(1\le p\le2)
\]
for a universal $K$.  Fixed-gap expanders then give the required lower
bounds, while the known $\ell_p$ upper bound transfers through the
diagonal isometry $\ell_p\hookrightarrow S_p$.
\end{abstract}

\section{Source question}

The source is Francisco Escudero Gutierrez and Garazi Muguruza,
\emph{All $S_p$ notions of quantum expansion are equivalent},
arXiv:2405.03517v1 (2024), Question 2 in Section 3.2, PDF pages 4--5
\cite{source}.  It asks:

\begin{quote}
Let $p\in[1,\infty)$. What is the minimum distortion $D_{n,S_p}$
needed, as a function of $n$ and $p$, to embed any $n$-point metric
space into $S_p$?
\end{quote}

The following are direct crops from the source PDF.  They include the
question, the definition of $\widetilde h_{S_p}$, and the proposed
expander route to a lower bound.

\clearpage
\begin{center}
\includegraphics[width=\textwidth]{figures/open_problem_crop_page4.png}

\vspace{0.5em}
\includegraphics[width=\textwidth]{figures/open_problem_crop_page5.png}
\end{center}

\section{Proof intuition}

There are two complementary regimes.  For $p\ge2$, the mixed function
space $L_p(V;S_p)$ is itself a noncommutative $L_p$ space.  Interpolating a
centered lazy random walk between its Hilbert spectral gap and its
$L_\infty$ bound makes a $k$-step walk contractive once
$k=O(p/\delta)$.  Jensen's inequality and a path decomposition then turn
this into the desired one-step Poincare inequality.

That interpolation necessarily degenerates at $p=1$.  For
$1\le p\le2$, translate the map at a vertex chosen by averaging and apply
the noncommutative Mazur map into the Hilbert space $S_2$.  Its inverse
estimate contains a radial weight $r^{2-p}$.  The chosen origin bounds the
total radial moment, and a weighted Holder inequality shows that at least
half of the original pairwise $p$-energy survives after the transformation.
The ordinary Hilbert Poincare inequality then gives a bound uniform all the
way to the trace-class endpoint.

\section{Definitions and results}

Let $G=(V,E)$ be a finite undirected $d$-regular graph, $n=|V|$,
and let $P$ be its normalized adjacency matrix.  Write
\[
 \delta=1-\lambda_2(P),
\]
where $\lambda_2(P)$ is the largest eigenvalue on the orthogonal
complement of the constants.  For $1\le p<\infty$, set
\[
 \widetilde h_{S_p}(G)=
 \inf_{f:V\to S_p}
 \frac{\left(\frac1{|E|}\sum_{\{i,j\}\in E}
       \|f(i)-f(j)\|_{p}^{p}\right)^{1/p}}
      {\left(\frac1{n^2}\sum_{i,j\in V}
       \|f(i)-f(j)\|_{p}^{p}\right)^{1/p}},
\]
where constant maps are omitted.

\begin{theorem}[Schatten-valued graph Poincare inequalities]\label{thm:graph}
There is a universal constant $K\ge1$ such that every finite regular
graph with $\delta>0$ satisfies
\[
 \widetilde h_{S_p}(G)\ge
 \begin{cases}
 (\delta/K)^{1/p},&1\le p\le2,\\[4pt]
 \delta/(14p),&2\le p<\infty.
 \end{cases}
\]
In particular, on any fixed-gap expander family the first bound is uniform
for $1\le p\le2$, and the second is $\Omega(1/p)$ for $p\ge2$.
\end{theorem}

The source proposes the stronger graph-by-graph comparison
$\widetilde h_{S_p}(G)\gtrsim_d\widetilde h_{S_2}(G)/p$.  Since
$\widetilde h_{S_2}(G)=\sqrt\delta$, our bounds have weaker dependence on
a gap tending to zero.  They have the precise uniformity needed on
fixed-gap expanders and hence suffice to answer the metric question.

\begin{corollary}[Full answer to Question 2]\label{cor:distortion}
There are universal constants $c,C>0$ such that for every $n\ge2$ and
every $1\le p<\infty$,
\[
 c\max\left\{1,\frac{\log n}{p}\right\}
 \le D_{n,S_p}\le
 C\max\left\{1,\frac{\log n}{p}\right\}.
\]
\end{corollary}

\section{Mazur-map lemmas for the endpoint range}

The lazy-walk interpolation proof for $p>1$ degenerates as $p\downarrow1$.
The replacement is nonlinear.  The Mazur map sends $S_p$ to the Hilbert
space $S_2$, but for $p<2$ its inverse Lipschitz estimate carries a
radial weight.  Choosing the translation point by averaging makes the
total radial moment no larger than the pairwise energy.  A second Holder
inequality then absorbs the weight without any loss depending on $p$.

For the polar decomposition $x=u|x|$, define
\[
 M_{p,2}(x)=u|x|^{p/2}.
\]
The identity $M_{2,p}(M_{p,2}(x))=x$ will transfer quantitative control
back from $S_2$ to $S_p$.

\begin{lemma}[Uniform Mazur estimates]\label{lem:mazur}
There is a universal constant $A\ge1$ such that, for every
$1\le p\le2$ and every $x,y\in S_p$,
\begin{align}
 \|M_{p,2}(x)-M_{p,2}(y)\|_2
 &\le A\|x-y\|_p^{p/2},\label{eq:mazur-forward}\\
 \|x-y\|_p
 &\le A\max\{\|x\|_p,\|y\|_p\}^{1-p/2}
       \|M_{p,2}(x)-M_{p,2}(y)\|_2.\label{eq:mazur-inverse}
\end{align}
The usual value $0^0=1$ may be used at $p=2$.
\end{lemma}

\begin{proof}
Ricard proved that the noncommutative Mazur map $M_{r,s}$ is
$\min\{r/s,1\}$-Holder on the unit ball of every noncommutative
$L_r$-space \cite{ricard}.  The quantitative proof gives a constant
depending only on the target exponent $s$ when $r<s$, and an absolute
constant times $r/s$ when $r>s$; see in particular the commutator
estimate in Lemma 2.6 of \cite{ricard}.  Thus the constants for
$M_{p,2}$, $1\le p\le2$, and for $M_{2,p}$, $1\le p\le2$, are
bounded by one universal constant.

The unit-ball estimate for $M_{p,2}$, together with its homogeneity of
degree $p/2$, gives \eqref{eq:mazur-forward} for arbitrary $x,y$.
For the inverse estimate, put $a=M_{p,2}(x)$, $b=M_{p,2}(y)$, and
$R=\max\{\|a\|_2,\|b\|_2\}$.  Homogeneity of $M_{2,p}$ and its
unit-ball Lipschitz bound give
\[
 \|M_{2,p}(a)-M_{2,p}(b)\|_p
 \le A R^{2/p-1}\|a-b\|_2.
\]
Since $\|a\|_2=\|x\|_p^{p/2}$ and $M_{2,p}(a)=x$, this is precisely
\eqref{eq:mazur-inverse}.  The zero case follows by continuity.  At
$p=2$, both maps are the identity.
\end{proof}

The next elementary estimate is the step that removes the radial weight.

\begin{lemma}[Weighted pair averaging]\label{lem:weighted}
Let $1\le p<2$, $q=2-p$, and let $x_i\in S_p$ for $i\in V$.
Write
\[
 d_{ij}=\|x_i-x_j\|_p,\qquad r_i=\|x_i\|_p,
 \qquad B=\E_{i,j}d_{ij}^p.
\]
If $\E_i r_i^p\le B$, then, with a zero summand when its denominator
vanishes,
\[
 \E_{i,j}\frac{d_{ij}^2}{r_i^q+r_j^q}\ge\frac B2.
\]
\end{lemma}

\begin{proof}
Put $s_{ij}=r_i^q+r_j^q$ and
$C=\E_{i,j}d_{ij}^2/s_{ij}$.  Holder's inequality with conjugate
exponents $2/p$ and $2/q$ gives
\begin{align*}
 B
 &=\E_{i,j}
   \left(\frac{d_{ij}^2}{s_{ij}}\right)^{p/2}s_{ij}^{p/2}\\
 &\le C^{p/2}\left(\E_{i,j}s_{ij}^{p/q}\right)^{q/2}.
\end{align*}
Because $p/q\ge1$,
\[
 (r_i^q+r_j^q)^{p/q}
 \le2^{p/q-1}(r_i^p+r_j^p).
\]
Consequently
\[
 \E_{i,j}s_{ij}^{p/q}
 \le2^{p/q}\E_i r_i^p\le2^{p/q}B.
\]
It follows that
\[
 B\le2^{p/2}C^{p/2}B^{q/2}.
\]
If $B>0$, division and raising to the power $2/p$ give $B\le2C$.
The case $B=0$ is immediate.
\end{proof}

\section{Proof of Theorem~\ref{thm:graph} for
\texorpdfstring{$1\le p\le2$}{1 <= p <= 2}}

Fix $f:V\to S_p$, and abbreviate
\[
 d_{ij}=\|f(i)-f(j)\|_p,
 \qquad B=\E_{i,j}d_{ij}^p.
\]
There is a vertex $o\in V$ such that
\[
 \E_i d_{io}^p\le B,
\]
because the average of the left side over $o$ equals $B$.  Translate
at this vertex: $x_i=f(i)-f(o)$, $r_i=\|x_i\|_p=d_{io}$, and set
\[
 a_i=M_{p,2}(x_i),\qquad h_{ij}=\|a_i-a_j\|_2.
\]

Suppose first that $p<2$, and put $q=2-p$.  Squaring
\eqref{eq:mazur-inverse} gives
\[
 h_{ij}^2
 \ge A^{-2}\frac{d_{ij}^2}{\max\{r_i,r_j\}^{q}}
 \ge A^{-2}\frac{d_{ij}^2}{r_i^q+r_j^q}.
\]
Lemma~\ref{lem:weighted} therefore implies
\begin{equation}\label{eq:pair-lower}
 \E_{i,j}h_{ij}^2\ge\frac{B}{2A^2}.
\end{equation}
On the other hand, \eqref{eq:mazur-forward} gives on every edge
\begin{equation}\label{eq:edge-upper}
 h_{ij}^2\le A^2d_{ij}^p.
\end{equation}

The Hilbert-space Poincare inequality for the normalized adjacency matrix
is
\begin{equation}\label{eq:hilbert-poincare}
 \E_{i,j}\|a_i-a_j\|_2^2
 \le\frac1\delta\frac1{|E|}
       \sum_{\{i,j\}\in E}\|a_i-a_j\|_2^2.
\end{equation}
Indeed, both sides equal twice the corresponding variance and Dirichlet
form normalization.  Combining \eqref{eq:pair-lower},
\eqref{eq:edge-upper}, and \eqref{eq:hilbert-poincare} yields
\[
 B\le\frac{2A^4}{\delta}\frac1{|E|}
       \sum_{\{i,j\}\in E}d_{ij}^p.
\]
Thus the first line of Theorem~\ref{thm:graph} holds with $K=2A^4$.
For $p=2$, it is just \eqref{eq:hilbert-poincare}, so increasing $K$
if necessary covers the closed interval $1\le p\le2$.  Notice that no
limiting argument is used at $p=1$.

\section{Proof of Theorem~\ref{thm:graph} for
\texorpdfstring{$p\ge2$}{p >= 2}}

For completeness, we include the independent interpolation argument that
has the correct $1/p$ dependence in the high-$p$ range.  Put
$Q=(I+P)/2$, and let $J$ be averaging over $V$.  All eigenvalues of
$Q$ lie in $[0,1]$, and on the orthogonal complement of constants,
\[
 \alpha:=\|Q-J\|_{\ell_2(V)\to\ell_2(V)}=1-\frac\delta2<1.
\]
Moreover, $QJ=JQ=J$, so $Q^k-J=(Q-J)^k$.

Consider the semifinite von Neumann algebra
\[
 \mathcal A=\ell_\infty(V)\,\overline\otimes\,B(\ell_2)
\]
with normalized trace in the $V$-coordinate.  Then
$L_r(\mathcal A)=L_r(V;S_r)$ isometrically.  The maps $Q^k$ and $J$
are contractions on $L_\infty(\mathcal A)$, whereas at the Hilbert
endpoint
\[
 \|Q^k-J\|_{L_2(\mathcal A)\to L_2(\mathcal A)}=\alpha^k.
\]
Noncommutative complex interpolation \cite{kosaki} gives, for $p\ge2$,
\begin{equation}\label{eq:interp-high}
 \|Q^k-J\|_{L_p(\mathcal A)\to L_p(\mathcal A)}
 \le2^{1-2/p}\alpha^{2k/p}.
\end{equation}

Let $L=-\log\alpha$, interpreting $\alpha=0$ by taking $k=1$, and
choose
\[
 k=\max\left\{1,
 \left\lceil\frac{(p-1)\log2}{L}\right\rceil\right\}.
\]
Then \eqref{eq:interp-high} is at most $1/2$.  For $f:V\to S_p$, set
$\bar f=Jf$ and $z=f-\bar f$.  Centering, the contraction, Jensen's
inequality, and Minkowski's inequality give
\begin{align}
 \left(\E_{i,j}\|f(i)-f(j)\|_p^p\right)^{1/p}
 &\le2\|z\|_{L_p(V;S_p)}\nonumber\\
 &\le4\left(\frac1n\sum_{i,j}(Q^k)_{ij}
       \|f(i)-f(j)\|_p^p\right)^{1/p}.\label{eq:kstep}
\end{align}

Let $(X_t)_{t=0}^k$ be the stationary lazy random walk with transition
matrix $Q$.  The triangle inequality and convexity imply
\[
 \|f(X_0)-f(X_k)\|_p^p
 \le k^{p-1}\sum_{t=1}^k
       \|f(X_{t-1})-f(X_t)\|_p^p.
\]
After averaging, stationarity shows that the $k$-step energy in
\eqref{eq:kstep} is at most $k^p$ times the one-step lazy energy.  Its
non-loop part is one half of the source edge energy.  Dropping this harmless
factor gives
\[
 \left(\E_{i,j}\|f(i)-f(j)\|_p^p\right)^{1/p}
 \le4k\left(\frac1{|E|}\sum_{\{i,j\}\in E}
       \|f(i)-f(j)\|_p^p\right)^{1/p}.
\]
Thus $\widetilde h_{S_p}(G)\ge1/(4k)$.  Finally,
$-\log(1-\delta/2)\ge\delta/2$ and $0<\delta\le2$, so
\[
 k\le\frac{(2+2\log2)p}{\delta}.
\]
Since $4(2+2\log2)<14$, this proves the second line of
Theorem~\ref{thm:graph}.

\section{Proof of Corollary~\ref{cor:distortion}}

Fix a constant-degree family of classical expanders with spectral gaps at
least $\delta_0>0$ \cite{hlw}.  In a $d$-regular $n$-vertex graph, a
ball of radius $r$ contains at most $1+d+\cdots+d^r$.  Hence a fixed
positive fraction of ordered vertex pairs have graph distance
$\Omega_d(\log n)$, and their $p$-root mean distance is
$\Omega_d(\log n)$, uniformly in $p\ge1$.

Suppose that the graph metric embeds in $S_p$ with distortion $D$, and
rescale the embedding to be nonexpanding.  The edge-energy root is at most
$1$, while the all-pairs energy root is at least $c_d\log n/D$.  When
$1\le p\le2$, Theorem~\ref{thm:graph} bounds the latter by a universal
constant because $(\delta_0/K)^{1/p}\ge\delta_0/K$.  Thus
\[
 D\gtrsim\log n\qquad(1\le p\le2).
\]
When $p\ge2$, the high-$p$ graph bound gives
\[
 D\gtrsim\frac{\log n}{p}.
\]
Together with $D\ge1$, these estimates are
$D\gtrsim\max\{1,\log n/p\}$ throughout the full range.  Standard
fixed-degree expanders at all sufficiently large admissible sizes, and the
monotonicity of worst-case distortion, give the conclusion for every $n$.

For the upper bound, Matousek's extension of Bourgain's construction gives
distortion $O(\log n/p)$ into $\ell_p$ in the range $p\lesssim\log n$
\cite{matousek}; the source paper recalls the same estimate \cite{source}.
The diagonal map $\ell_p\hookrightarrow S_p$ is isometric.  When
$p\ge\log n$, the Frechet embedding into $\ell_\infty^n$, followed by
the identity into $\ell_p^n$, has distortion at most
$n^{1/p}\le e$.  Enlarging the universal constant handles the finitely
many small cases and proves the matching upper bound.

\section{Verification, limitations, and novelty}

The weighted averaging computation and finite matrix instances of the
Mazur-energy argument were checked by
\texttt{code/check\_constants.py}; these checks are consistency tests, not
proof.  A separate normalization and endpoint audit appears in
\texttt{verification.md}.  The original source PDF and Ricard's supporting
paper are included beside this packet.

The conclusion is full relative to Question 2: it determines the sharp
joint asymptotics in $n$ and $p$ for every $1\le p<\infty$.  It does
not prove the stronger graph-by-graph comparison suggested in the source;
for a small spectral gap our displayed dependence on $\delta$ is weaker.

A bounded novelty search was performed on 19 July 2026.  The run indexes
were searched using arXiv:2405.03517, its exact title, $D_{n,S_p}$,
Schatten distortion, noncommutative Mazur maps, trace norm, and graph
Poincare keywords.  Targeted arXiv/web searches used the exact source
question and close phrases.  Ricard's quantitative Mazur-map paper
(arXiv:1407.8334), Mendel--Naor's superexpander paper (arXiv:1207.4705),
and related Schatten nonembedding results were checked.  No later paper
explicitly settling Question 2, including its $S_1$ endpoint, was found.
Novelty confidence is moderate because this was a bounded search rather
than an exhaustive citation review.

The most important human-review point is the uniformity in $p\in[1,2]$
of Lemma~\ref{lem:mazur}.  It follows from the constants in Ricard's proof,
not merely from the qualitative statement of his theorem.  The remaining
new ingredient, Lemma~\ref{lem:weighted}, is proved above with the explicit
constant $1/2$.

\begin{thebibliography}{9}

\bibitem{source}
F. Escudero Gutierrez and G. Muguruza,
\emph{All $S_p$ notions of quantum expansion are equivalent},
arXiv:2405.03517v1, 2024.

\bibitem{hlw}
S. Hoory, N. Linial, and A. Wigderson,
\emph{Expander graphs and their applications},
Bull. Amer. Math. Soc. 43 (2006), 439--561.

\bibitem{kosaki}
H. Kosaki,
\emph{Applications of the complex interpolation method to a von Neumann
algebra: non-commutative $L^p$-spaces},
J. Funct. Anal. 56 (1984), 29--78.

\bibitem{matousek}
J. Matousek,
\emph{On embedding expanders into $\ell_p$ spaces},
Israel J. Math. 102 (1997), 189--197.

\bibitem{ricard}
E. Ricard,
\emph{Holder estimates for the noncommutative Mazur maps},
Arch. Math. (Basel) 104 (2015), 37--45; arXiv:1407.8334.

\end{thebibliography}

\end{document}
